\documentclass[10pt,a4paper]{amsart}
\usepackage[margin=19mm]{geometry}
\usepackage{amsmath,amssymb,mathtools}
\usepackage[scr=rsfs,cal=euler]{mathalfa}
\usepackage{xfrac}
\usepackage[unicode,hidelinks,bookmarks=false]{hyperref}
\newcommand{\ie}{i.e.\ }
\newcommand{\cf}{cf.\ }
\newcommand{\resp}{respectively\ }
\newcommand{\Subsection}{\subsection}
\providecommand{\nobreakdash}{\nobreak\hbox{-}}
\setlength{\emergencystretch}{2em}
\multlinegap0pt
\usepackage{amssymb}
\usepackage[shortlabels]{enumitem}
\newtheorem{theorem}{Theorem}
\newtheorem{corollary}{Corollary}
\title{Preliminaries on Pre-Hilbert Structures on Polynomial Spaces and Associated Laplacians}
\author{Jean-Pierre Magnot}
\date{Source: arXiv:2512.22394v2 (CC BY 4.0)}
\begin{document}
\maketitle

\noindent\emph{Source and licence.} Preliminaries on Pre-Hilbert Structures on Polynomial Spaces and Associated Laplacians. Version arXiv:2512.22394v2 (CC BY 4.0), distributed by arXiv under the Creative Commons Attribution 4.0 International licence, http://creativecommons.org/licenses/by/4.0/, as recorded in the arXiv licence metadata for this exact version and shown on the abstract page https://arxiv.org/abs/2512.22394v2. Full licence notice: \texttt{licenses/arxiv-2512.22394-CC-BY-4.0.txt}.\par\smallskip

\noindent\textit{Editorial scope.} Counted unit: the article's corollary cor:kernel-stability (source0.tex lines 1308--1324). The article's complete original proof of it is delivered with it (source0.tex lines 1326--1429, one \texttt{proof} environment of 104 lines). No mathematics is written, repaired or re-derived: the statement and the proof are the article's own lines, in the article's own notation, and no retained line is substituted or re-flowed.  Retained uncounted as the source-local context the proof needs: lines 1223--1238.  Nothing was omitted from any retained span, so the package carries no omission record.  No retained line contains a citation, so the wrapper carries no bibliography and no bibliography entry.  No retained line contains a figure environment, an image-inclusion command or drawing code, so no image is delivered and the package declares no asset dependency.  Editorial formatting only, in addition: an AMS \texttt{amsart} preamble that restates the article's own declarations and macros used by the retained lines, namely the environments \texttt{theorem}, \texttt{corollary} on separate counters, together with the standard packages those lines need.  The retained environments print in this document as Theorem 1, Corollary 1 in order of appearance, whereas the article prints the same items with the numbering of its own full document.  The complete native source of the article as distributed in the arXiv e-print for this exact version is bundled unchanged as \texttt{sources/191732/source0.tex}.
\begin{theorem}[Stability of polynomial projectors]
\label{thm:projector-stability}
Let $\mathfrak G_1,\mathfrak G_2$ be two polynomial Hilbert geometries acting on the same
Hilbert space $H$, and assume
\[
d_{\mathrm{res}}(\mathfrak G_1,\mathfrak G_2)
=
\big\|(1+\Delta_1)^{-1}-(1+\Delta_2)^{-1}\big\|
\le \varepsilon.
\]
Then, for every fixed degree $N$, there exists a constant $C_N>0$ such that
\[
\|P_N^{(1)}-P_N^{(2)}\|_{\mathcal B(H)}
\le C_N\,\varepsilon.
\]
\end{theorem}

\begin{corollary}[Kernel stability on compacts]\label{cor:kernel-stability}
Assume $\Omega$ is a metric space and fix a compact set $K\subset\Omega$.
For $i=1,2$, let
\[
K_N^{(i)}(x,y)=\sum_{k=0}^N p_k^{(i)}(x)\overline{p_k^{(i)}(y)}
\]
be the polynomial reproducing kernel of $H_{\le N}$ (in any orthonormal basis).
Then there exists a constant $C_{N,K}>0$ such that
\[
\sup_{(x,y)\in K\times K}\big|K_N^{(1)}(x,y)-K_N^{(2)}(x,y)\big|
\le C_{N,K}\,\|P_N^{(1)}-P_N^{(2)}\|_{\mathcal B(H)}.
\]
In particular, under the assumptions of Theorem~\ref{thm:projector-stability}, one has
\[
\sup_{K\times K}|K_N^{(1)}-K_N^{(2)}|\le C'_{N,K}\,\varepsilon.
\]
\end{corollary}

\begin{proof}
Let $V_i:=\mathrm{Ran}(P_N^{(i)})\subset H$ (finite-dimensional), and set
\[
V:=V_1+V_2\subset H.
\]
All vectors in $V$ are (restrictions of) polynomial functions on $\Omega$, hence are continuous on $K$.
For each $x\in K$, consider the conjugate evaluation functional on $V$,
\[
\Lambda_x:V\to\mathbb C,\qquad \Lambda_x(f):=\overline{f(x)}.
\]
Since $V$ is finite-dimensional, $\Lambda_x$ is continuous for the $H$-norm.
By the Riesz representation theorem (for continuous conjugate-linear functionals under our convention),
there exists a unique vector $g_x\in V$ such that
\begin{equation}\label{eq:Riesz-gx}
\Lambda_x(f)=\langle f,g_x\rangle_H\qquad\forall f\in V.
\end{equation}

\smallskip
\noindent\textbf{Step 1: reproducing vectors in $V_i$.}
For $i\in\{1,2\}$ define
\[
k_x^{(i)}:=P_N^{(i)}g_x\in V_i.
\]
Then for every $f\in V_i$ we have, using $P_N^{(i)}f=f$ and orthogonality of $P_N^{(i)}$,
\[
\overline{f(x)}=\Lambda_x(f)=\langle f,g_x\rangle_H=\langle P_N^{(i)}f,g_x\rangle_H
=\langle f,P_N^{(i)}g_x\rangle_H=\langle f,k_x^{(i)}\rangle_H.
\]
Thus $k_x^{(i)}$ is the Riesz representer of $\Lambda_x$ restricted to $V_i$.

\smallskip
\noindent\textbf{Step 2: identification with the kernel.}
Let $\{p_0^{(i)},\dots,p_N^{(i)}\}$ be any orthonormal basis of $V_i$ in $H$.
Expanding $k_x^{(i)}$ in this basis and using $\langle p_j^{(i)},k_x^{(i)}\rangle_H=\overline{p_j^{(i)}(x)}$,
we obtain
\begin{equation}\label{eq:kx-expansion}
k_x^{(i)}=\sum_{j=0}^N \overline{p_j^{(i)}(x)}\,p_j^{(i)}.
\end{equation}
Therefore, for $x,y\in K$,
\[
\langle k_y^{(i)},k_x^{(i)}\rangle_H
=\sum_{j=0}^N p_j^{(i)}(y)\,\overline{p_j^{(i)}(x)}
=K_N^{(i)}(y,x).
\]
In particular,
\begin{equation}\label{eq:kernel-via-k}
\big|K_N^{(1)}(x,y)-K_N^{(2)}(x,y)\big|
=\big|K_N^{(1)}(y,x)-K_N^{(2)}(y,x)\big|
=\big|\langle k_y^{(1)},k_x^{(1)}\rangle_H-\langle k_y^{(2)},k_x^{(2)}\rangle_H\big|.
\end{equation}

\smallskip
\noindent\textbf{Step 3: estimate by $\|P_N^{(1)}-P_N^{(2)}\|$.}
Using $k_x^{(i)}=P_N^{(i)}g_x$ and adding/subtracting,
\[
\langle P_1 g_y,P_1 g_x\rangle_H-\langle P_2 g_y,P_2 g_x\rangle_H
=
\langle (P_1-P_2)g_y,P_1 g_x\rangle_H+\langle P_2 g_y,(P_1-P_2)g_x\rangle_H,
\]
where we wrote $P_i:=P_N^{(i)}$ for brevity.
Hence, since $\|P_i\|=1$,
\[
\big|\langle k_y^{(1)},k_x^{(1)}\rangle_H-\langle k_y^{(2)},k_x^{(2)}\rangle_H\big|
\le
\|P_1-P_2\|_{\mathcal B(H)}\big(\|g_y\|_H\,\|g_x\|_H+\|g_y\|_H\,\|g_x\|_H\big)
=
2\|P_1-P_2\|\,\|g_x\|\,\|g_y\|.
\]
Taking the supremum over $(x,y)\in K\times K$ and using \eqref{eq:kernel-via-k}, we obtain
\begin{equation}\label{eq:kernel-bound}
\sup_{(x,y)\in K\times K}\big|K_N^{(1)}(x,y)-K_N^{(2)}(x,y)\big|
\le
2\Big(\sup_{x\in K}\|g_x\|_H\Big)^2\,\|P_N^{(1)}-P_N^{(2)}\|_{\mathcal B(H)}.
\end{equation}

\smallskip
\noindent\textbf{Step 4: finiteness of the constant on compacts.}
It remains to check that
\[
M_{N,K}:=\sup_{x\in K}\|g_x\|_H<\infty.
\]
Choose a basis $(\phi_1,\dots,\phi_m)$ of $V$ and let $G\in\mathbb C^{m\times m}$ be its Gram matrix
$G_{\alpha\beta}:=\langle \phi_\beta,\phi_\alpha\rangle_H$ (invertible since it is positive definite on $V$).
Writing $g_x=\sum_{\beta} c_\beta(x)\phi_\beta$ and using \eqref{eq:Riesz-gx} with $f=\phi_\alpha$ gives
\[
\overline{\phi_\alpha(x)}=\langle \phi_\alpha,g_x\rangle_H=\sum_{\beta} c_\beta(x)\langle \phi_\alpha,\phi_\beta\rangle_H
=\sum_{\beta} G_{\alpha\beta}\,c_\beta(x),
\]
hence $c(x)=G^{-1}\overline{\Phi(x)}$ with $\Phi(x)=(\phi_1(x),\dots,\phi_m(x))^\top$.
Since each $\phi_\alpha$ is continuous on $K$, the map $x\mapsto \Phi(x)$ is continuous on the compact set $K$,
so $\sup_{x\in K}\|\Phi(x)\|<\infty$. Therefore $\sup_{x\in K}\|c(x)\|<\infty$, and hence
$\sup_{x\in K}\|g_x\|_H<\infty$. This proves $M_{N,K}<\infty$ and thus \eqref{eq:kernel-bound}.

Setting $C_{N,K}:=2M_{N,K}^2$ yields the first inequality in the corollary.

\smallskip
\noindent\textbf{Final step.}
Under the assumptions of Theorem~\ref{thm:projector-stability}, we have
$\|P_N^{(1)}-P_N^{(2)}\|_{\mathcal B(H)}\le C_N\,\varepsilon$, hence
\[
\sup_{K\times K}|K_N^{(1)}-K_N^{(2)}|\le (C_{N,K}C_N)\,\varepsilon,
\]
so one may take $C'_{N,K}:=C_{N,K}C_N$.
\end{proof}

\end{document}
